\section{Appendix: Scope and Qualifications}
This study concerns the Twenty-first Sunday after Pentecost in the 1962
Roman Missal on 18 October 2026 under the general calendar. No diocesan,
religious, or titular calendar has been supplied. The mission addition
and St Luke commemoration have the conditions stated in the opening;
neither imports another Mass's readings or Preface. The two readings
are editorial constructions from compatible arguments in the cited
witnesses. Their authors did not all preach on this complete formulary.

The principal early witnesses are Chrysostom, Augustine, and Gregory.
Aquinas supplies a later theological distinction; Bellarmine supplies
later Catholic Psalm exposition; Schuster offers a modern historical
liturgical reading. Their differences of purpose and period remain
visible. Chrysostom's disputed explanations of Ephesians are identified
against his historical English editor's notes. Augustine's and
Chrysostom's strong warnings about forfeited pardon are not erased by
Aquinas's more precise distinction. Augustine's future judgment and
Schuster's present deprivation of truth are related interpretations,
not identical glosses. Contextual ancient social examples and polemic
are not silently promoted into pastoral prescriptions.

The study depends on the acquired historical editions and retained
transcriptions listed below. It does not offer a fresh critical Greek
or Latin edition. The lawful Missal transcription and historical English
are taken from the reviewed text collation; normalized typography is
not a new translation. Modern source introductions and the current
Missal scan are evidentiary sources with their own rights, not a license
to reproduce their complete publications. Exact machine-readable source
bindings and evidence limits remain in the canonical research dossier.

The work's substantive boundaries are deliberate: no new local calendar
decision, textual conjecture, or claim of exhaustive scholarly consensus
is supplied. The study asks how these appointed words can be read
coherently in the company of identifiable Christian interpreters,
while leaving their actual reasoning and disagreements available to
the reader.

\section{References}
\subsection{Liturgical and biblical witnesses}
\begin{itemize}
\item \work{Missale Romanum}, Vatican typical edition (1962): Sunday
pp.~413--415, nos.~1699--1708; Trinity Preface p.~293, no.~1082;
St Luke pp.~689--690, nos.~3883,3889,3891; propagation of the faith
pp.~[86]--[87], nos.~4833,4841,4843. Controlling
\href{https://media.churchmusicassociation.org/pdf/missale62.pdf}{CMAA facsimile}.
General Rubrics 91,95,107--111; General Rubrics of the Missal
444--446,453,480,483,494,500--501,505,509 govern the dated resolution.
\item Public-domain Latin antecedents: Pustet \work{Missale Romanum}
(1862), Sunday pp.~354--355; Mame Tours \work{Missale Romanum}
(1922), St Luke pp.~738--739.
\item \work{The Roman Missal}, Eugene Cummiskey (Philadelphia, 1861):
Sunday pp.~450--452; Trinity Preface pp.~xxviii--xxix; St Luke
pp.~677--678. Historical human English, with abbreviated conclusions.
\item F.~X. Lasance and F.~A. Walsh, \work{The New Roman Missal},
revised 1945 printing: propagation of the faith pp.~1453,1456,1457.
The acquired file's URL says 1937; the inspected printing is revised
1945. Only the identified inherited prayer bodies are used here.
\item \work{Douay--Rheims Bible}, Challoner revision,
\href{https://www.gutenberg.org/ebooks/1581}{Project Gutenberg 1581}:
Esther, Job, Psalms, Matthew, Ephesians, and the biblical echoes
identified in the commentary. Vulgate Psalm numbering is retained.
\end{itemize}

\subsection{Exegetical witnesses actually used}
\begin{itemize}
\item John Chrysostom, \work{Homilies on Matthew} LXI, \S\S1--5,
NPNF first series 10, historical English revised by George Prevost
and M.~B. Riddle, retained CCEL transcription. The complete homily
frames the unmerciful servant through received gifts and refused mercy.
\item John Chrysostom, \work{Homilies on Ephesians} XXII--XXIV,
NPNF first series 13, William John Copeland's historical English,
revised by Gross Alexander; retained CCEL text. Consult XXII on
Eph.~6:9--13, XXIII on 6:14, XXIV on 6:14--17, including editorial
notes 465 and 468 at the disputed expressions.
\item Augustine, \work{Sermon} 83, NPNF first series 6,
\work{Sermons on New Testament Lessons} XXXIII, \S\S1--8,
R.~G. MacMullen's historical English, retained CCEL text.
\item Augustine, \work{Expositions of the Psalms}, NPNF first series 8,
retained CCEL historical English: Ps.~89 \S\S1--4 (NPNF Ps.~90);
Ps.~113 \S\S1--6 (NPNF Ps.~114); Ps.~118 Aleph, verses~1--8,
and Caph, verses~81--88 (NPNF Ps.~119). Section and octonary
locators disambiguate the differing Psalm numbers.
\item Augustine, \work{On the Trinity} III, preface \S3, Arthur West
Haddan's historical English (1887), retained
\href{https://www.newadvent.org/fathers/130103.htm}{New Advent delivery}.
The delivery identifies later editorial revision; the argument about
divine missions is paraphrased here, not quoted as untouched 1887 bytes.
\item Gregory the Great, \work{Moralia in Iob} II.10.16--11.19 and
III.3.3--6.8, \work{Library of the Fathers} English (1844), retained
complete Lectionary Central transcription. The reset transcription,
not a newly inspected printed facsimile, is the authoring witness.
\item Thomas Aquinas, \work{Summa theologiae} III, q.~88, aa.~1--4,
English Dominican translation, retained
\href{https://www.newadvent.org/summa/4088.htm}{New Advent text}:
especially a.~1, response and reply~1, and a.~2, response.
\item Robert Bellarmine, \work{A Commentary on the Book of Psalms},
John O'Sullivan's translation (1866), retained ecatholic2000 text:
Ps.~89:1--2;113:1;118:81,84,86,88. Later exposition, not patristic
authorship of these psalms.
\item Alfredo Ildefonso Schuster, \work{The Sacramentary}, trans.
Arthur Levelis-Marke, III (London: Burns Oates \& Washbourne, 1927),
pp.~178--181, ``Twenty-first Sunday after Pentecost,'' retained
Internet Archive OCR with the acquired volume facsimile. The argument
here uses the text and its specific prayer comments, not unchecked
claims about earlier chant forms.
\end{itemize}

Historical date sources are cited individually in the preceding appendix.
